\documentclass[UTF8,11pt]{ctexart}
\usepackage[a4paper,margin=22mm]{geometry}
\usepackage{graphicx,enumitem,amssymb,tabularx}
\usepackage[normalem]{ulem}
\usepackage{xeCJKfntef}
\setlength{\parindent}{0pt}
\setlength{\parskip}{.6em}
\sloppy
\begin{document}
\section*{2023年6月英语四级真题(第3套)}
\subsection*{2023 年6 月大学英语四级考试真题（第3 套）}


\subsection*{\textbf{Part I Writing} \textbf{( 30 minutes)}}


\textbf{Directions:}\textit{ Suppose your university is conducting a survey to collect students’ opinions of online classes. You} \textit{are to write a response to the survey about their advantages and disadvantages, and what improvements can} \textit{be made. You will have 30 minutes for the task. You should write at least }\textit{\uline{120 }}\textit{words but no more than }\textit{\uline{180}} \textit{words.}


\subsection*{\textbf{Part II} \textbf{Listening Comprehension} \textbf{( 25 minutes)}}


提示：2023 年6 月全国只考两套听力，故本套未重复显示。


\subsection*{\textbf{Part III Reading Comprehension ( 40 minutes)}}


\subsection*{\textbf{Section A}}


\textbf{Directions:} \textit{In this section, there is a passage with ten blanks. You are required to select one word for each} \textit{blank from a list of choices given in a word bank following the passage. Read the passage through carefully} \textit{before making your choices. Each choice in the bank is identified by a letter. Please mark the corresponding} \textit{letter for each item on }\textbf{\textit{Answer Sheet 2}}\textit{ with a single line through the centre. You may not use any of the words} \textit{in the bank more than once.}


Morocco is responding to increasing energy demands by setting up one of the largest solar plants in the world.


The Noor solar power station is \uline{26} in the city of Ouarzazate and, once completed, will generate 580 million watts of electricity. The World Bank estimates it will serve 1.1 million people. It’s \uline{27} to be completed soon.


Morocco’s current energy comes \uline{28} from imports. The nation hopes to get 50 percent of its energy from renewable sources by 2030. With demand for energy \uline{29} at an annual rate of 7 percent, the new solar plant could be a \uline{30} part of that goal.


“This makes Morocco a big \uline{31} in the field of solar energy in the Arab region and the African continent. It could also be a forerunner for many other countries in the world that \uline{32} on foreign imports of energy,” said Ali Hajji, a solar energy specialist and engineering professor.


Experts believe that the Middle East and North Africa have huge \uline{33} for solar energy projects. This is partly because of adequate sunlight and partly because technology has become more \uline{34} in the region.


“The last few years have seen a realization of \uline{35} how competitive solar technologies can be,” said Michael Taylor, a senior analyst at the International Renewable Energy Agency.


\begin{itemize}[leftmargin=2em,itemsep=.2em,topsep=.4em]\item[A.] affordable
\item[B.] ancestor
\item[C.] crucial
\item[D.] depend
\item[E.] initial
\item[F.] insist
\item[G.] just
\item[H.] located
\item[I.] mostly
\item[J.] operating
\item[K.] perhaps
\item[L.] pioneer
\item[M.] potential
\item[N.] rising
\item[O.] scheduled
\end{itemize}

\subsection*{\textbf{Section B}}


\textbf{Directions:} \textit{In this section, you are going to read a passage with ten statements attached to it. Each statement} \textit{contains information given in one of the paragraphs. Identify the paragraph from which the information is} \textit{derived. You may choose a paragraph more than once. Each paragraph is marked with a letter. Answer the} \textit{questions by marking the corresponding letter on }\textbf{\textit{Answer Sheet 2}}\textit{.}


\textbf{New Formula One Chief Hopes to Grab Americans’ Attention}


A) For the past four decades, the leader of Formula One car racing, one of the biggest annual sporting series in the world, was Bernie Ecclestone, a former motorcycle parts dealer who built it into an international presence essentially on his own.


B) A skilled backroom operator who speaks without a filter, Ecclestone said often that in his opinion, the sport was at its best when he was allowed to act as “a dictator.”


C) Yet now the dictator is gone. After an American company, Liberty Media, acquired the Formula One competition recently, Chase Carey — a former executive with Fox Broadcasting Company and DirecTV who by his own admission is not a fierce racing fan — was named to replace Ecclestone and to try to renovate the organization’s management, reach and ambition.


D) Among the goals, Carey said in an interview on Tuesday, is one that just about every global sport seems interested in chasing: increasing interest in the United States. “People have said we’re going to ‘Americanize’ it,” Carey said. “And we’re not going to do that totally. But realistically, there are some elements of Americanization that the sport could use.”


E) While Formula One commands enormous audiences throughout much of the world, many American sports fans know it as that other motorsport, the one that is not Nascar (纳斯卡车赛). Formula One teams race far more technologically advanced vehicles around tracks all over the world — in magnificent events in places like Malaysia, Monaco, Singapore and the United Arab Emirates, and on tradition-rich tracks like Silverstone in England and Monza in Italy too.


F) The series has an annual race in Austin, Texas. But within “a few years,” Carey said, he plans to bring another to a destination American city, like New York, Los Angeles, Miami or Las Vegas. Carey’s ambitious plan is two-fold: first, change the business model of Formula One, which he said was a “one-man show” under Ecclestone that had a largely narrow vision when it came to negotiating partnership deals; and second, alter the way fans experience the sport, both in person and remotely, so that connections between the audience and people within the series are easier to make.


G) Increased digital access for fans, a more behind-the-scenes experience for broadcast viewers and innovation in areas like virtual reality — what is it like to speed around a track inside a Ferrari? — are among the possibilities. “The sport has clearly been underserved,” Carey said. “It doesn’t do anything digitally. There’s no marketing. It doesn’t tell any stories. The goal in this is to make the fans connect to the live experience as much as possible, and the tools you have to do that, we’re not using at all.”


H) The larger question, though, is a familiar one: Is there room for Formula One in the ever-crowded sports landscape of the United States? Opinions vary, particularly because viewing habits among consumers continue to evolve. John Bloom, a professor at Shippensburg University who has studied American sports history, said the biggest challenge for any sport trying to increase its presence in the United States was framing itself in a way that had lasting appeal. “Sports generally become popular in some way because they establish a narrative,” Bloom said. “When I think of motorsports in the U.S., what we all think of is Nascar, and the narrative of Nascar is sort of rural, white, working-class Americans, mostly in the South, connecting with the atmosphere of those races. That’s the narrative. When I think of the narrative of Formula One, it’s a very different kind of audience.”


I) That difference, Carey said, is significant. While some might immediately link Formula One to Nascar in terms of American growth, Carey said Formula One’s brand research had indicated there was very little crossover; rather, Formula One fans generally cite other so-called elite events, like Wimbledon or the Ryder Cup, as competitions they enjoy. “Other than they’re both cars, the Nascar fan base is a very different fan base,” Carey said. “It’s a very regional fan base. Formula One is a global, famous brand of stars. These are machines that shock and awe you.”


J) Carey’s background is in deal making and innovation. At Fox Broadcasting Company, he was a top advisor for years, known for his skill in helping to lead the launch of company into sports, as well as the start of Fox News Channel. After going to DirecTV, he positioned the satellite provider as a mainstream option in millions of households.


K) Now, after Liberty Media paid \$4.4 billion to acquire Formula One, he is charged with making the investment pay off. “I think they can build Formula One in the U.S.,” said Patrick Crakes, an executive at InVivo Media Group who spent 25 years at Fox Broadcasting Company before leaving in 2016 as a senior manager at Fox Sports. “People don’t work on their cars anymore. They don’t want that connection anymore. It’s about technology and pushing the limits. It’s about speed, danger and risk. And Formula One has that more than any other racing series.”


L) That is what hooked Carey, and he said he thought his experience was not unusual. He recalled attending Formula One’s Monaco race last year and being overwhelmed by the ceremony leading up to the event, the way the race charmed the city for days ahead of the start. In his mind it felt like a Super Bowl (超级碗橄榄球赛).


M) Then, on race day, he watched as the cars rocketed out of a tunnel and went screaming toward a tight turn with the city’s harbor and the Mediterranean Sea in the background framing the scene. He was fascinated. “You can’t help but be awed,” he said, “and I think that feeling can be translated to the viewer.”


N) He added: “The broader sport is a little too inward-looking, and we need to be more open. In some ways, I’m glad to be coming from the outside. The guys who are in the sport forever are sitting there saying: ‘We can’t do that. We can’ t do that because it’s never been done that way.’ ”


36. Chase Carey believes greater use should be made of digital technology to make Formula One more


accessible to its fans.


37. Chase Carey was deeply impressed by the ceremony preceding last year’s Monaco race.


38. One of Chase Carey’s goals is to make Formula One more appealing to Americans.


39. A former motorbike parts dealer led Formula One for the past forty years.


40. Chase Carey thought the audience of Formula One could be made to share his feeling about the race.


41. Chase Carey used to serve as a top advisor for a major broadcasting company.


42. Chase Carey intends to make connections easier between the audience and the Formula One racers.


43. The new leader of Formula One admitted he was not super interested in car racing.


44. People’s opinions differ as to whether Formula One can be promoted in the U.S.


45. Compared with other racing series, Formula One focuses more on speed and involves more danger.


\subsection*{\textbf{Section C}}


\textbf{Directions: }\textit{There are 2 passages in this section. Each passage is followed by some questions or unfinished} \textit{statements. For each of them there are four choices marked A), B), C) and D). You should decide on the best} \textit{choice and mark the corresponding letter on }\textbf{\textit{Answer Sheet 2}}\textit{ with a single line through the centre.}


\subsection*{\textbf{Passage One}}


\textbf{Questions 46 to 50 are based on the following passage.}


Supermarkets have long been suffering as one of the thinnest-margined businesses in existence and one of the least-looked-forward-to places to work or visit. For more than a decade, they have been under attack from e-commerce giants, blamed for making Americans fat, and accused of contributing to climate change.


Supermarkets can technically be defined as giants housing 15,000 to 60,000 different products. The revolutionary idea of a self-service grocery, where people could hunt and gather food from aisles rather than asking a clerk to fetch items from behind a counter, first came about in America. There is some debate about which was the very first, but over the years a consensus has built around King Kullen Supermarket, founded in New York in 1930.


For some 300 years, Americans had fed themselves from small stores and public markets. Shopping for food involved mud, noisy chickens, clouds of flies, nasty smells, bargaining, and getting short-changed. The supermarket imitated the Fordist factory, with its emphasis on efficiency and standardization, and reimagined it as a place to buy food. Supermarkets may not feel cutting-edge now, but they were a revolution in distribution at the time. They were such strange marvels that, on her first official state visit to the United States in 1957, Queen Elizabeth II insisted on an impromptu (即兴的) tour of a suburban-Maryland Giant Food.


The typical supermarket layout has barely changed over the past 90 years. Most stores open with flowers, fruit and vegetables at the front as a breath of freshness to arouse our appetite. Meanwhile, they keep the milk, eggs, and other daily basics all the way back so you’ll travel through as much of the store as possible, and be tempted along the way. In the early days, as the supermarket multiplied, so did our suspicion of it. We have long feared that this “revolution in distribution” uses corporate black magic on our appetite. The book \textit{The Hidden} \textit{Persuaders}, published in 1957, warned that supermarkets were putting women in a “hypnoidal trance (催眠恍惚状态),” causing them to wander aisles bumping into boxes and “picking things off shelves at random.”


46. What problem have supermarkets been facing?


\begin{itemize}[leftmargin=2em,itemsep=.2em,topsep=.4em]\item[A.] They are actually on the way to bankruptcy.
\item[B.] They have been losing customers and profits.
\item[C.] They are forced to use e-commerce strategies.
\item[D.] They have difficulty adapting to climate change.
\end{itemize}

47. What does the passage say about the idea of a self-service grocery?


\begin{itemize}[leftmargin=2em,itemsep=.2em,topsep=.4em]\item[A.] It was put forward by King Kullen.
\item[B.] It originated in the United States.
\item[C.] It has been under constant debate.
\item[D.] It proves revolutionary even today.
\end{itemize}

48. What did supermarkets do by adopting the Fordist factory approach?


\begin{itemize}[leftmargin=2em,itemsep=.2em,topsep=.4em]\item[A.] They modernized traditional groceries in many ways.
\item[B.] They introduced cutting-edge layout of their stores.
\item[C.] They improved the quality of the food they sold.
\item[D.] They revolutionized the distribution of goods.
\end{itemize}

49. What is the typical supermarket layout intended to do?


\begin{itemize}[leftmargin=2em,itemsep=.2em,topsep=.4em]\item[A.] Arouse customers’ appetite to buy flowers, fruit and vegetables.
\item[B.] Provide customers easy access to items they want to buy.
\item[C.] Induce customers to make more unplanned purchases.
\item[D.] Enable customers to have a more enjoyable shopping experience.
\end{itemize}

50. What have people long feared about supermarkets?


\begin{itemize}[leftmargin=2em,itemsep=.2em,topsep=.4em]\item[A.] They use tricky strategies to promote their business.
\item[B.] They are going to replace the local groceries entirely.
\item[C.] They apply corporate black magic to the goods on display.
\item[D.] They take advantage of the weaknesses of women shoppers.
\end{itemize}

\subsection*{\textbf{Passage Two}}


\textbf{Questions 51 to 55 are based on the following passage.}


The traditional school year, with three months of vacation every summer, was first implemented when America was an agricultural society and the summer months were needed for farm work. Since then, we’ve completely changed as a nation. Students no longer spend summers farming, but they aren’t in school, either. The average American student receives 13 weeks off from school each calendar year — with about 11 of those during the summer. Few other countries have more than seven weeks off in a school calendar.


With the US lagging behind countries in academics, it’s time to consider year-round schooling. One benefit of this change is that students will not fall victim to the “summer slide,” or the well-documented phenomenon where students forget some of the knowledge they have acquired when too much time is taken off from school. Decades of research shows that it can take from 8 to 13 weeks at the beginning of every school year for students to get back to where they were before the summer holiday.


But year-round schooling isn’t just about academics. Teachers and students experience a closer relationship in year-round schools than they do in traditional schools and, in the absence of any long-term break, students do not feel detached from the school environment. These closer bonds and greater attachment pay off. Research shows that students in year-round schools are more self-confident and feel more positive about their schooling experience.


But don’t kids need time to relax? Some childhood development experts believe that time off from school is vital to healthy development as kids are not designed to spend so much of their time inside classrooms and the summer break provides a perfect opportunity to get outside. The problem with this argument is that most children aren’t playing outside or even spending time with other kids. While some children visit summer camps, most stay at home, watching TV or playing games on electronic devices, which hardly benefits them.


The US has changed from a farming economy to a knowledge- and innovation-based economy, so it makes sense for the school year to change as well.


51. Why did America’s traditional school year have a three-month summer vacation?


\begin{itemize}[leftmargin=2em,itemsep=.2em,topsep=.4em]\item[A.] Students needed to help with farm work.
\item[B.] Students needed time to learn necessary farming skills.
\item[C.] The agricultural society then attached less importance to academics.
\item[D.] America lagged behind other countries in making a scientific school calendar.
\end{itemize}

52. What benefit will year-round schooling bring students in addition to improving their learning?


\begin{itemize}[leftmargin=2em,itemsep=.2em,topsep=.4em]\item[A.] It will help them get back to where their lessons started.
\item[B.] It will enable them to absorb what they have learned.
\item[C.] It will familiarize them with the school environment.
\item[D.] It will strengthen their relationship with teachers.
\end{itemize}

53. What do some childhood development experts believe about the long summer vacation?


\begin{itemize}[leftmargin=2em,itemsep=.2em,topsep=.4em]\item[A.] It meets students’ need to study on their own.
\item[B.] It enables students to learn about the outside world.
\item[C.] It satisfies students’ desire to stay longer at home.
\item[D.] It contributes to students healthy growth.
\end{itemize}

54. What is the argument against the experts’ idea of a long summer vacation?


\begin{itemize}[leftmargin=2em,itemsep=.2em,topsep=.4em]\item[A.] It does little good to most students.
\item[B.] It benefits few students playing outside.
\item[C.] It leads students to neglect their studies.
\item[D.] It makes students addicted to computer games.
\end{itemize}

55. What does the author think of the traditional school year in the US today?


\begin{itemize}[leftmargin=2em,itemsep=.2em,topsep=.4em]\item[A.] Well-grounded.
\item[B.] Culture-bound.
\item[C.] Outdated.
\item[D.] Welcomed.
\end{itemize}

\subsection*{\textbf{Part IV Translation} \textbf{( 30 minutes)}}


\textbf{Directions: }\textit{For this part, you are allowed 30 minutes to translate a passage from Chinese into English. You} \textit{should write your answer on }\textbf{\textit{Answer Sheet 2}}\textit{.}


中国越来越重视终身教育，发展继续教育是构建终身教育体系的有效途径。高校作为人才培养的基地，拥有先进的教学理念和优越的教学资源，理应成为继续教育的办学主体。因此，近年来许多高校适应社会需求，加强与用人单位沟通，努力探索一条符合中国国情的继续教育发展新路，以使继续教育在国家发展战略中发挥更大的作用。

\end{document}